\documentclass[UTF8]{ctexart}
\usepackage{geometry}
\usepackage{hyperref}
\usepackage{listings}
\usepackage{xcolor}

\geometry{a4paper, left=2.5cm, right=2.5cm, top=2.5cm, bottom=2.5cm}

\title{Agent开发周报 260201}
\author{刘德辰}
\date{2026年02月01日}

\begin{document}

\maketitle

\section{交信投测试环境适配与接口标准化}

本周重点完成了 Agent 系统与交信投测试环境的接口适配工作，实现了“即插即用”的架构设计。

\begin{itemize}
    \item \textbf{API Client 封装}：在 \texttt{agent/src/api-client.ts} 中实现了标准化的 \texttt{APIClient} 类。该客户端严格遵循交信投接口文档规范，封装了交信投提供的具体单任务数据接口，包括预警清单 (\texttt{getAlerts})、车辆实时工况 (\texttt{getVehicleRealtime})、热力图数据 (\texttt{getHeatmap}) 及驾驶员画像 (\texttt{getDriverFeatures}) 等核心数据接口。这些接口均为具体的 REST API 端点，每个接口负责单一的数据获取任务。
    \item \textbf{接口规范一致性}：尽管 Demo 界面目前仍使用自有的 Demo 数据或后端服务，但其请求结构、参数定义及响应格式已与交信投标准接口文档完全对齐。
    \item \textbf{无缝切换}：通过统一的配置管理，系统仅需切换 Base URL 和 API Key 即可从 Demo 环境无缝迁移至交信投真实测试环境，无需修改业务逻辑代码。
\end{itemize}

\section{工作场景识别重构与用户自维护}

用工作场景列表的形式重构了路由拒绝逻辑，该方案支持动态管理和用户自定义。

\begin{itemize}
    \item \textbf{D1 数据库存储}：将原有的硬编码场景迁移至 D1 数据库 (\texttt{work\_scenarios} 表)，定义了包括 \texttt{name} (场景名)、\texttt{description} (LLM 语义描述)、\texttt{tool} (绑定工具)、\texttt{required\_params} (必填参数) 等字段。
    \item \textbf{双重匹配机制}：Router 模块升级为“向量检索 + LLM 判决”的双重匹配机制。系统首先通过 Embedding 向量相似度筛选 Top-K 候选场景，再由 Router LLM 最终判断用户意图是否命中特定场景，有效提高了意图识别的准确率。
    \item \textbf{用户维护界面}：实现了工作场景的管理功能，允许用户（管理员）在前端界面动态新增、修改或禁用工作场景。这使得 Agent 的能力边界可以随业务需求灵活调整，而无需重新部署代码。
\end{itemize}

\section{Demo 前端功能优化}

针对 Demo 前端（\texttt{frtend-tsx}）进行了多项功能增强，以配合 Agent 能力的展示。

\begin{itemize}
    \item \textbf{场景管理 UI}：新增了工作场景配置页面，支持可视化编辑场景描述和参数要求。
    \item \textbf{交互体验}：优化了 Agent 对话界面的响应流程，适配了流式输出（Streaming）和 Markdown 渲染，提升了用户与 Omni 智能体交互的流畅度。
    \item 增加了会话重命名和删除功能。
\end{itemize}

\section{通用智能体体验优化}

优化了通用智能体（Omni Agent）的交互逻辑，使其更贴近实际业务使用场景。

\begin{itemize}
    \item \textbf{冷启动体验提升}：针对用户初次进入或发送简单的“你好”等问候语时，智能体不再输出机械式的通用介绍，而是根据当前角色（公交安全助手）提供简洁、亲切的引导，直接进入服务状态。
    \item \textbf{聚焦业务场景}：通过 Prompt 调优，限制了智能体在非业务相关话题上的发散，确保对话始终围绕公交安全运营、辅助决策等核心领域展开，避免提及无关的通用场景信息。
\end{itemize}

\section{工具层逻辑抽象化}

对底层工具架构进行了重构，实现了逻辑与数据的解耦。

\begin{itemize}
    \item \textbf{ToolProvider 抽象层}：引入了 \texttt{ToolProvider} 接口，将数据获取逻辑与上层智能体技能（Skill）彻底解耦。当前实现为 \texttt{LocalToolProvider}，通过统一的 \texttt{query\_data} 工具从 D1 数据库读取数据。智能体不再直接依赖特定的数据库实现，而是通过抽象的 \texttt{ToolProvider} 接口获取数据。该设计为未来接入 MCP（Model Context Protocol）协议预留了扩展空间，只需切换 \texttt{ToolProvider} 实现即可，无需修改业务逻辑代码。这不仅简化了 \texttt{generate\_driver\_report} 等工具的实现逻辑，也提高了代码的可扩展性和可维护性。
\end{itemize}

\end{document}
